\documentclass[11pt]{article}
\usepackage[margin=1in]{geometry}
\usepackage{amsmath}
\usepackage{amssymb}
\usepackage{enumitem}
\usepackage{hyperref}
\usepackage{listings}
\usepackage{xcolor}

\title{CSCE 355: Turing Machine Assignment 2}
\author{}
\date{}

\lstset{
  basicstyle=\ttfamily\small,
  breaklines=true,
  frame=single,
  backgroundcolor=\color{gray!10}
}

\begin{document}
\maketitle

% ------------------------------------------------------------------
% NOTE TO STUDENTS: every machine you submit MUST be deterministic.
% A nondeterministic machine receives NO CREDIT for that problem,
% even if it happens to give correct accept/reject answers on the
% test strings. See Section "Determinism Requirement" below for
% exactly how to check this in JFLAP before you submit.
% ------------------------------------------------------------------

\section*{Overview}

In this assignment you will build five deterministic single-tape Turing
machines in JFLAP, one per problem below. Each machine decides (accepts or
rejects) a language -- your machine does not need to leave any particular
output on the tape, it only needs to halt in an accepting state on strings
that belong to the language, and halt in a non-accepting state (or fail to
halt) on strings that do not.

You will submit five \texttt{.jff} files. Each will be tested automatically
against 20 provided test strings (10 that your machine must accept, 10 that
it must reject), and each problem is worth 5 points, for 25 points total.

\subsection*{Getting JFLAP}

This assignment uses JFLAP, a free tool for building and simulating
automata. Download it directly from the official site:

\begin{center}
\url{https://www.jflap.org}
\end{center}

JFLAP is a \texttt{.jar} file and requires a working Java installation
(Java 8 or later is recommended). It is not included in this assignment
package -- download it yourself from the link above. Do not use a
third-party mirror.

\section*{Determinism Requirement}

\textbf{Every machine you submit must be deterministic.} If a machine is
found to be nondeterministic, it receives \textbf{zero credit} for that
problem, regardless of whether it produces correct accept/reject results
on the test strings. This is checked automatically by the grading script
before any test strings are even run.

A Turing machine is deterministic if, for every state and every tape
symbol, there is \emph{at most one} transition that applies. Concretely,
watch out for these two situations:

\begin{enumerate}[label=(\alph*)]
  \item \textbf{Two transitions leaving the same state on the same read
        symbol.} If state $q_i$ has two different outgoing transitions
        that both read the symbol \texttt{0} (even if they write different
        symbols, move in different directions, or go to different states),
        the machine is nondeterministic at $q_i$.
  \item \textbf{A transition that can apply simultaneously with another.}
        In a standard JFLAP Turing machine, every transition reads some
        tape symbol (including the blank symbol, shown as \texttt{$\square$}), so
        this usually reduces to case (a) above -- but double check any
        state with several outgoing arrows to make sure no two of them
        share a read symbol.
\end{enumerate}

\subsection*{How to check for determinism in JFLAP}

\textbf{Method 1: Test $\to$ Highlight Nondeterminism.} With your machine
open in the editor, use the menu option \textbf{Test $\to$ Highlight
Nondeterminism}. JFLAP will open a new view of your machine with any
nondeterministic states highlighted. If nothing is highlighted, your
machine passes this check. If any state is highlighted, inspect its
outgoing transitions and remove or fix the conflicting one(s).

\textbf{Method 2: Watch the number of simultaneous configurations while
stepping.} A deterministic machine only ever has \emph{one} active
configuration at any point during simulation. When you step through your
machine (\textbf{Input $\to$ Step}), if you ever see more than one
configuration displayed at the same time -- or if JFLAP asks you to
freeze/remove configurations because there are too many active at once to
display -- that is a clear sign your machine has taken a nondeterministic
branch somewhere. A correctly built deterministic machine for this
assignment should never show more than one configuration box during
stepping.

Fix any state flagged by either method before submitting. The autograder
performs an equivalent check on the raw transition table of your
\texttt{.jff} file, so a machine that fails determinism in JFLAP will
also fail it (and receive 0 points for that problem) when graded.

\section*{Problems}

Each problem below is worth 5 points. Alphabet symbols other than those
listed (and the blank symbol) should not appear in valid input, and your
machine does not need to handle them specially.

\begin{enumerate}

\item \textbf{Palindrome.} Input alphabet $\{0,1\}$. Accept if and only if
the input string is a palindrome (reads the same forwards and backwards).
The empty string counts as a palindrome.
\\ Files: \texttt{problem1\_palindrome.jff}, \texttt{problem1\_palindrome\_tests.txt}

\item \textbf{Unary Addition Check.} Input alphabet $\{0, \#, =\}$. Input
strings are of the form $0^m \# 0^n = 0^p$ (unary numbers $m$, $n$, $p$
separated by \texttt{\#} and \texttt{=}). Accept if and only if $m + n =
p$.
\\ Files: \texttt{problem2\_addition.jff}, \texttt{problem2\_addition\_tests.txt}

\item \textbf{Binary Increment Check.} Input alphabet $\{0, 1, \#\}$.
Input strings are of the form $x \# y$ where $x$ and $y$ are binary
numbers. Accept if and only if $y$ equals $x + 1$.
\\ Files: \texttt{problem3\_increment.jff}, \texttt{problem3\_increment\_tests.txt}

\item \textbf{Copy Check.} Input alphabet $\{0, 1, \#\}$. Input strings
are of the form $w \# w'$ where $w, w' \in \{0,1\}^*$. Accept if and only
if $w = w'$ (the two halves are identical).
\\ Files: \texttt{problem4\_copy.jff}, \texttt{problem4\_copy\_tests.txt}

\item \textbf{$0^n1^n2^n$.} Input alphabet $\{0, 1, 2\}$. Accept if and
only if the input has the form $0^n 1^n 2^n$ for some $n \geq 0$ (some
number of 0's, followed by the same number of 1's, followed by the same
number of 2's).
\\ Files: \texttt{problem5\_anbncn.jff}, \texttt{problem5\_anbncn\_tests.txt}

\end{enumerate}

\section*{Test String Files}

Each \texttt{*\_tests.txt} file contains exactly 20 lines:

\begin{itemize}
  \item \textbf{Lines 1--10} are strings your machine must \textbf{accept}.
  \item \textbf{Lines 11--20} are strings your machine must \textbf{reject}.
\end{itemize}

A blank line represents the empty string ($\lambda$). You can load a
test file directly into JFLAP's batch simulator via
\textbf{Input $\to$ Multiple Run}, then \textbf{File $\to$ Open} the
\texttt{*\_tests.txt} file from within that window to load all 20 strings
at once, and click \textbf{Run Inputs} to see accept/reject results for
each row.

\textbf{Do not edit the \texttt{*\_tests.txt} files.} They will be
replaced with the original versions before grading, so any changes you
make will simply be discarded -- editing them will only make it harder
for you to debug locally with an inaccurate picture of what will actually
be graded.

\section*{Testing Your Machines}

You have two complementary ways to test your work as you go:

\begin{enumerate}
\item \textbf{Inside JFLAP (interactive).} Use \textbf{Input $\to$
Step} to walk through individual strings one transition at a time and
watch the tape, or \textbf{Input $\to$ Fast Run} for a quick single-string
accept/reject check, or \textbf{Input $\to$ Multiple Run} loaded with the
provided \texttt{*\_tests.txt} file to check all 20 strings for a problem
at once.

\item \textbf{From the command line (automated).} This assignment folder
includes \texttt{test\_tm.py}, the same grading script that will be used
to grade your submission. From inside this folder, run:
\begin{lstlisting}
python3 test_tm.py
\end{lstlisting}
This will parse all five of your \texttt{.jff} files, check each for
determinism, run all 20 test strings per problem, and print a full report
along with your total score out of 25. Run this as often as you like
while working -- it is exactly what will be run when your assignment is
graded (with the test files and this script itself refreshed to their
original versions first, so make sure you are testing against the
originals).
\end{enumerate}

\section*{Grading}

Each of the five problems is worth 5 points, for 25 points total. For a
given problem, you receive full credit (5 points) only if \emph{both} of
the following hold:

\begin{itemize}
  \item Your submitted machine is deterministic (see above), \textbf{and}
  \item Your machine produces the correct accept/reject result on all 20
        test strings for that problem.
\end{itemize}

Otherwise, that problem receives 0 points. There is no partial credit
within a single problem.

\section*{Submission Instructions}

\begin{enumerate}
  \item Open \textbf{YOURUSERNAME.txt} and put your university username
        as the \emph{first line} of the file (and only your username --
        nothing else on that line). This is how the grading script
        identifies your submission.
  \item Build and save your five Turing machines into the provided
        \texttt{.jff} files (\texttt{problem1\_palindrome.jff} through
        \texttt{problem5\_anbncn.jff}), overwriting the empty templates.
  \item Verify all five machines are deterministic and pass their test
        strings using \texttt{python3 test\_tm.py} as described above.
  \item Submit this entire folder (or a zip of it) as instructed by your
        TA/instructor. Do not rename any of the provided files.
\end{enumerate}

\end{document}
